\chapter{Example: the gate-driven isolated dot}\label{ch:ex-dot}

The first example is an exactly solvable, isolated normal dot. It
introduces three concepts used throughout the following chapters: quasienergies, harmonic weights
(photon sidebands) and the Floquet zone.

\section{Exact solution}

Consider the isolated impurity \eqref{eq:Hd} with a gate drive,
\begin{equation}
  H(t)=\Hd(\eps_0)+A\cos(\Omega t)\,(n_d-1).
\end{equation}
Every Fock state $\ket n$ of the dot, with charge $n\in\{0,1,2\}$, is an eigenstate at
every time, so its time evolution consists only of a phase factor,
\begin{equation}
  \ket{\psi_n(t)}=\ee^{-\ii E_nt}\,\ee^{-\ii q_n(A/\Omega)\sin\Omega t}\ket n,\qquad q_n=n-1 .
\end{equation}
Comparing with \cref{eq:floquet-state}, the quasienergy is the static energy, $\qe_n=E_n$,
for every drive amplitude: this purely diagonal, zero-mean modulation preserves the
time-averaged energy. The Floquet mode is the periodic phase factor. With the generating
function of the Bessel functions,
$\ee^{-\ii z\sin\vartheta}=\sum_mJ_m(z)\ee^{-\ii m\vartheta}$,
\begin{equation}\label{eq:dot-weights}
  \ket{u_n(t)}=\sum_m \ee^{\ii m\Omega t}(-1)^mJ_m\!\left(\frac{q_nA}\Omega\right)\ket n,
  \qquad \norm{u_{n,m}}^2=J_m^2\!\left(\frac{q_nA}{\Omega}\right).
\end{equation}
The harmonic weights depend only on the \emph{modulation index} $q_nA/\Omega$, the ratio
between the amplitude of the level oscillation and the photon energy. \Cref{fig:dot}(a)
compares them with the Sambe results.

\begin{physics}
$J_m^2(A/\Omega)$ is the probability that a driven level has exchanged $m$ photons with
the field. An electron that tunnels from a lead into such a level can do so at any of the
energies $E_n+m\Omega$, so the tunnelling spectrum acquires sidebands with weights
$J_m^2$. This is the Tien--Gordon theory of photon-assisted tunnelling, first formulated for
superconducting tunnel junctions~\cite{tien1963multiphoton,dayem1962quantum} and observed in
quantum dots~\cite{kouwenhoven1994photon,kouwenhoven1994observation,oosterkamp1998microwave}.
For a modulation index of $2.405$, the first zero of $J_0$, the central line disappears.
\end{physics}

\section{Gauge and the meaning of the charge}

Only charge \emph{differences} enter a physical transition: tunnelling from $\ket n$ to
$\ket{n+1}$ involves the relative modulation index $(q_{n+1}-q_n)A/\Omega=A/\Omega$. The
absolute weights \eqref{eq:dot-weights} depend on the constant in $\Hd$. Writing the drive as
$A\cos(\Omega t)\,n_d$ instead would assign $q_n=n$, and the empty dot would have no
sidebands at all. The two forms differ by the c-number $A\cos\Omega t$, a gauge
transformation (\cref{sec:gauge}) that leaves the quasienergies unchanged. \pkg{} keeps the
constants of the model by default, so that $q_n=n-1$ is measured from half filling; this
centres the sidebands and minimizes the Floquet cutoff needed for convergence. The option
\code{drop_scalar_harmonics=True} removes the oscillating c-number terms, which here gives
the second form.

\section{Floquet zones and replicas}

\Cref{fig:dot}(b) shows the replicas $\qe_n+m\Omega$ of the empty and doubly occupied
states as a function of $\Omega$ at fixed amplitude, with line intensity proportional to
$J_m^2(A/\Omega)$. All replicas are equally valid quasienergies; the shaded first zone
$\abs\qe<\Omega/2$ is one choice of representatives. At low frequency the modulation index
is large, many sidebands carry weight and the replicas fan out; at high frequency the
weight concentrates in $m=0$. Replicas of different states cross freely: without a coupling
between the states, crossings in the Floquet spectrum are exact.

\begin{figure}[htb]
\centering
\manualfigure{dot}
\caption{Gate-driven isolated dot, $U=\q{dot.u}$, $\eps_0=\q{dot.epsilon}$. (a) Harmonic weights
of the doubly occupied state versus the modulation index $A/\Omega$; circles are Sambe results, lines
the Bessel functions \eqref{eq:dot-weights}. (b) Replicas $\qe+m\Omega$ of $\ket0$ and
$\ket2$ versus $\Omega$ for $A=\q{dot.fan_amplitude}$; line opacity and width indicate the harmonic
weight. The shaded region is the first Floquet zone.}
\label{fig:dot}
\end{figure}

\section{In the code}

\begin{lstlisting}
import numpy as np
from qdjj_solver import Hamiltonian, Impurity, PeriodicHamiltonian, Sector
from qdjj_solver.qp_solver import solve_floquet

def dot(detuning):
    impurity = Impurity.anderson(1., detuning)
    return Hamiltonian(impurity.operator, 2, (1, -1), observables=impurity.observables)

A, omega = .6, 1.
periodic = PeriodicHamiltonian.from_callable(lambda t: dot(.2 + A*np.cos(omega*t)), omega,
                                             assume_fixed_coordinates=True)
states = solve_floquet(periodic, 0, Sector(0, 0), dict(floquet_cutoff=12, states=2))
print(states.quasienergies)            # 0.3 and 0.7: the static energies of |0> and |2>
print(states.harmonic_weights[:, states.harmonic_indices == 0])   # J_0(A/omega)^2 for both
\end{lstlisting}

The isolated dot is constructed directly in fixed physical-electron coordinates.
The explicit assertion \code{assume_fixed_coordinates=True} (\cref{sec:fixed-coordinates})
allows the code to compare the samples through their mode labels. Since the only harmonics of this drive
are $H_0$ and $H_{\pm1}$, any number of samples $N\ge3$ reproduces $H(t)$ exactly.

\tutorial{02_driven_dot}
